\section{Introduction}
Every workday morning the southbound lanes of National Freeway~1 between the Zhubei and Hsinchu
interchanges fill with commuters heading to Hsinchu Science Park; the 4.8\,km stretch goes from
about 3\,min at free flow to a median of \res{profile.workdayPeakMedian}\,min at
\res{profile.workdayPeakSlot}. Roadside signs show only current conditions. A commuter at home who
knows the conditions \emph{one hour ahead} can decide to leave now or to wait. We ask:

\emph{Using only information available at prediction time, can a linear model with upstream and
calendar features forecast the segment travel time 60 minutes ahead accurately enough---peak MAE
at most \res{cfg.threshold} minutes---to support a commuter's departure decision?}

We state four hypotheses (H1 persistence, H2 upstream propagation, H3 calendar, H4 rain at peak)
and test them with Ridge regression under a leakage-free, time-ordered protocol. Contributions:
(i)~workday$\times$\allowbreak slot interactions, motivated by the different \emph{shape} of workday
and non-workday profiles, are the most valuable feature group (Section~\ref{sec:results});
(ii)~upstream travel time adds information beyond the segment's own state
(Section~\ref{sec:analysis}); (iii)~the linear model beats gradient boosting in the peak but not
overall, and in routine peaks barely beats a historical average, so its value lies in
non-routine conditions.
